\documentclass[11pt]{article}

\usepackage[margin=1in]{geometry}
\usepackage[T1]{fontenc}
\usepackage{lmodern}
\usepackage{amsmath,amssymb,amsthm}
\usepackage{stmaryrd}
\usepackage{mathpartir}
\usepackage{booktabs}
\usepackage{placeins}

% let figures share pages with text, and keep them in the section that discusses them
\renewcommand{\topfraction}{0.9}
\renewcommand{\bottomfraction}{0.9}
\renewcommand{\textfraction}{0.05}
\renewcommand{\floatpagefraction}{0.8}
\setcounter{topnumber}{3}
\setcounter{bottomnumber}{3}
\setcounter{totalnumber}{5}
\usepackage[hidelinks]{hyperref}

% ---------- keywords ----------
\newcommand{\kw}[1]{\mathsf{#1}}
\newcommand{\str}{\kw{str}}
\newcommand{\lit}[1]{\kw{lit}\,#1}
\newcommand{\prop}{\kw{prop}}
\newcommand{\prf}[1]{\kw{prf}\;#1}
\newcommand{\pred}[1]{\kw{pred}\;#1}
\newcommand{\true}{\kw{true}}
\newcommand{\letin}[3]{\kw{let}\;#1 = #2\;\kw{in}\;#3}

% ---------- kinds ----------
\newcommand{\proper}{\ast}
\newcommand{\populated}{\ast^{+}}
\newcommand{\kall}[3]{\forall #1 : #2.\,#3}

% ---------- types ----------
\newcommand{\dict}[3]{\{[#1 : #2] : #3\}}
\newcommand{\gdict}[4]{\{[#1 : #2]\ #3 : #4\}}
\newcommand{\refine}[3]{\{#1 : #2 \mid #3\}}
\newcommand{\iabs}[3]{\lambda #1 : #2.\,#3}
\newcommand{\tabs}[3]{\lambda #1 : #2.\,#3}
\newcommand{\union}[1]{\textstyle\bigcup #1}
\newcommand{\iapp}[2]{#1\,[#2]}

% ---------- judgments ----------
\newcommand{\emptyctx}{\diamond}
\newcommand{\proves}{\Vdash}
\newcommand{\isprop}{\ \kw{prop}}
\newcommand{\discharge}[3]{#1 \vdash #2 \leadsto #3}
\newcommand{\kinding}[4]{#1 \mid #2 \vdash #3 : #4}
\newcommand{\sem}[1]{\llbracket #1 \rrbracket}
\newcommand{\rulename}[1]{\textsc{#1}}

% ---------- MiniMUMPS ----------
\newcommand{\false}{\kw{false}}
\newcommand{\ifte}[3]{\kw{if}\;#1\;\kw{then}\;#2\;\kw{else}\;#3}
\newcommand{\isdef}[1]{\kw{isdefined}\;#1}
\newcommand{\readm}[1]{\kw{read}\;#1}
\newcommand{\firstten}[3]{\kw{first10}\;#1\;\kw{with}\;#2\;\kw{in}\;#3}

% ---------- compilation to MiniMUMPS ----------
\newcommand{\dproc}[2]{\langle #1;\, #2 \rangle}
\newcommand{\supplies}[4]{#1;\, #2 \vdash #3 \Uparrow #4}
\newcommand{\matches}[4]{#1 \vdash #2 : #3 \rhd #4}
\newcommand{\compiles}[5]{#1;\, #2;\, #3 \vdash #4 \Rightarrow #5}
\newcommand{\sorttest}[2]{\lceil #1 : #2 \rceil}
\newcommand{\code}[1]{\texttt{#1}}

% the heading of a group of rules; mathpar already leaves space above the rules
\newcommand{\judgmentbox}[2]{\medskip\noindent\fbox{$#1$}\quad #2\par\vspace{-1.2\baselineskip}}

\emergencystretch=3em

\theoremstyle{remark}
\newtheorem*{remark}{Remark}

\title{The Surface, Core, and MiniMUMPS Languages of SchemaTypes}
\date{}

\begin{document}
\maketitle

\section{Overview}

SchemaTypes has two schema languages. Schemas are written in the \emph{surface} language. The
\emph{core} language is a conventional dependently sorted calculus, and every surface schema can be
translated to it. A surface schema can also be compiled to \emph{MiniMUMPS}, a small calculus for
reading a database, to obtain a program that tests whether a database conforms to the schema.

Both schema languages have two levels. The \emph{index} level describes the logical layer of a database:
strings, predicates on strings, and proofs that predicates hold. The \emph{type} level describes
the physical layer. The two languages differ mainly at the index level:
\begin{itemize}
\item In the core language, index-level functions are curried. A predicate is applied to one argument
  at a time, and its sort may interleave string arguments with proof arguments.
\item In the surface language, predicates are $n$-ary. A predicate is always applied to all of its
  string arguments at once, its proof arguments are never written, and its sort lists its string
  formals and the proofs it requires in two separate sections.
\end{itemize}
At the type level, both languages are curried: index-to-type and type-to-type functions take one
argument at a time. The type levels differ in one respect. A surface dictionary lists the
propositions that guard its keys, while the core language expresses each guard as a union over the
proofs of the guard.

This document defines the syntax of both schema languages (Sections~\ref{sec:surface-syntax}
and~\ref{sec:core-syntax}), the judgments that classify surface terms
(Section~\ref{sec:surface-judgments}), and the translation from the surface language to the core
language (Section~\ref{sec:translation}). It then defines MiniMUMPS
(Section~\ref{sec:minimumps}) and the compilation of surface schemas to MiniMUMPS
(Section~\ref{sec:compilation}). The two translations are independent of each other: compilation to
MiniMUMPS starts from the surface language, not from the core language.

The document does not define the judgments of the core language, the evaluation of types in either
schema language, or the denotational semantics of types. Section~\ref{sec:implementation} relates
the definitions to the F\# implementation.

\paragraph{Conventions.}
We write $a, b, P, Q$ for index variables, $X, Y$ for type variables, $s, t$ for string literals,
and $\ell$ for record labels. Index variables form a single set; we use $P$ and $Q$ for variables
that are meant to stand for predicates. A sequence $z_1, \ldots, z_n$ is abbreviated $\overline{z}$.
Terms are identified up to renaming of bound variables, and we assume that bound variables are
chosen so that the variables bound by a context are distinct. We write $[e/a]$ for the
capture-avoiding substitution of $e$ for $a$, as in $\kappa[e/a]$, and
$[\overline{e}/\overline{a}]$ for the simultaneous substitution of each $e_i$ for $a_i$.

\section{Surface syntax}
\label{sec:surface-syntax}

Figure~\ref{fig:surface-syntax} gives the abstract syntax of the surface language.

\begin{figure}[!htbp]
\[
\begin{array}{llcll}
\text{atom} & e & ::= & a & \text{index variable} \\
  & & \mid & s & \text{string literal} \\[4pt]
\text{predicate application} & \alpha & ::= & P(e_1, \ldots, e_n) & n \geq 0 \\[4pt]
\text{formula} & \varphi, \psi & ::= & \true & \\
  & & \mid & \alpha & \\[4pt]
\text{string sort} & c, d & ::= & \str & \text{all strings} \\
  & & \mid & \lit{s} & \text{only the string } s \\[4pt]
\text{signature} & S & ::= & (a_1 : c_1, \ldots, a_n : c_n \mid \alpha_1, \ldots, \alpha_m)
  & n, m \geq 0 \\[4pt]
\text{sort} & q & ::= & c & \text{strings} \\
  & & \mid & \prf{\varphi} & \text{proofs of } \varphi \\
  & & \mid & \pred{S} & \text{predicates} \\[4pt]
\text{kind} & \kappa, \rho & ::= & \proper & \text{proper types} \\
  & & \mid & \populated & \text{populated proper types} \\
  & & \mid & \kappa \to \rho & \text{type-to-type operators} \\
  & & \mid & \kall{a}{q}{\kappa} & \text{index-to-type operators} \\[4pt]
\text{type} & \tau, \sigma & ::= & X & \text{type variable} \\
  & & \mid & \gdict{a}{c}{\varphi_1, \ldots, \varphi_k}{\tau} & \text{dictionary, } k \geq 0 \\
  & & \mid & \{\ell_1 : \tau_1, \ldots, \ell_n : \tau_n\} & \text{record, } n \geq 1 \\
  & & \mid & \refine{a}{c}{\varphi} & \text{string refinement} \\
  & & \mid & \iabs{a}{q}{\tau} & \text{index-to-type abstraction} \\
  & & \mid & \tabs{X}{\kappa}{\tau} & \text{type-to-type abstraction} \\
  & & \mid & \union{\tau} & \text{union} \\
  & & \mid & \tau\;\sigma & \text{type application} \\
  & & \mid & \iapp{\tau}{e} & \text{index application} \\
  & & \mid & \letin{X}{\tau}{\sigma} & \text{type definition} \\[4pt]
\text{sort context} & \Omega, \Psi & ::= & \emptyctx \;\mid\; \Omega, a : q & \\[2pt]
\text{kind context} & \Gamma & ::= & \emptyctx \;\mid\; \Gamma, X : \kappa &
\end{array}
\]
\caption{Abstract syntax of the surface language}
\label{fig:surface-syntax}
\end{figure}

\paragraph{Index level.}
An \emph{atom} is an index variable or a string literal. Atoms are the only indices that can be
passed as arguments. A \emph{predicate application} $P(e_1, \ldots, e_n)$ applies the predicate
variable $P$ to all of its string arguments. A \emph{formula} is a predicate application or the
trivial proposition $\true$. Substitution also replaces a variable in predicate position:
$P(\overline{e})[Q/P] = Q(\overline{e})$.

A \emph{signature} $(a_1 : c_1, \ldots, a_n : c_n \mid \alpha_1, \ldots, \alpha_m)$ describes an
$n$-ary predicate. Its \emph{string section} $a_1 : c_1, \ldots, a_n : c_n$ names the predicate's
string formals and gives their sorts. Its \emph{proof section} $\alpha_1, \ldots, \alpha_m$ lists
propositions that must be proven wherever the predicate is applied. The formals
$a_1, \ldots, a_n$ are bound in the proof section. For example, a predicate with signature
\[ (\mathit{cust} : \str,\ \mathit{card} : \str \mid \mathit{CustId}(\mathit{cust})) \]
relates two strings, and may only be applied to $\mathit{cust}$ and $\mathit{card}$ where
$\mathit{CustId}(\mathit{cust})$ is known to hold.

There are three forms of sort. A string sort $c$ classifies strings, $\prf{\varphi}$ classifies
proofs of $\varphi$, and $\pred{S}$ classifies predicates with signature $S$. There is no sort of
propositions: formulas are a separate syntactic category, so no variable stands for a
proposition.

\paragraph{Type level.}
The keys of the dictionary type $\gdict{a}{c}{\varphi_1, \ldots, \varphi_k}{\tau}$ are the strings
$a$ of sort $c$ that satisfy the \emph{guards} $\varphi_1, \ldots, \varphi_k$, and its values have
type $\tau$. The key variable $a$ is bound in the guards and in $\tau$. A dictionary may have no
guards. In the
refinement $\refine{a}{c}{\varphi}$ the variable $a$ is bound in $\varphi$. In the kind
$\kall{a}{q}{\kappa}$ and the abstraction $\iabs{a}{q}{\tau}$ the variable $a$ is bound in $\kappa$
and $\tau$ respectively, but not in $q$. Record labels are distinct. A union $\union{\tau}$ is taken
over a type $\tau$ that is an index-to-type operator.

\FloatBarrier
\subsection{Concrete syntax and derived forms}
\label{sec:concrete}

Table~\ref{tab:concrete} shows how the abstract syntax is written. In the concrete syntax, index
variables, type variables, and record labels are identifiers, and string literals are quoted.

\begin{table}[!htbp]
\centering
\begin{tabular}{ll}
\toprule
Abstract & Concrete \\
\midrule
$P(e_1, \ldots, e_n)$ & \code{P($e_1$, \ldots, $e_n$)} \\
$\lit{s}$ & \code{"$s$"} \\
$\prf{\varphi}$ & \code{prf $\varphi$} \\
$\pred{(\overline{a : c} \mid \overline{\alpha})}$
  & \code{pred ($a_1$ :\ $c_1$, \ldots, $a_n$ :\ $c_n$ | $\alpha_1$, \ldots, $\alpha_m$)} \\
$\proper$, \ $\populated$ & \code{*}, \ \code{**} \\
$\kappa \to \rho$ & \code{$\kappa$ -> $\rho$} \\
$\kall{a}{q}{\kappa}$ & \code{($a$ :\ $q$) -> $\kappa$} \\
$\gdict{a}{c}{\varphi_1, \ldots, \varphi_k}{\tau}$, \ $k \geq 1$
  & \code{\{ [$a$ :\ $c$] $\varphi_1$, \ldots, $\varphi_k$ :\ $\tau$ \}} \\
$\{\ell_1 : \tau_1, \ldots, \ell_n : \tau_n\}$ & \code{\{ $\ell_1$ :\ $\tau_1$, \ldots, $\ell_n$ :\ $\tau_n$ \}} \\
$\refine{a}{c}{\varphi}$ & \code{\{ $a$ :\ $c$ | $\varphi$ \}} \\
$\iabs{a}{q}{\tau}$ & \code{($a$ :\ $q$) => $\tau$} \\
$\tabs{X}{\kappa}{\tau}$ & \code{($X$ ::\ $\kappa$) => $\tau$} \\
$\union{\tau}$ & see below \\
$\tau\;\sigma$ & \code{($\tau$ !$\sigma$)} \\
$\iapp{\tau}{e}$ & \code{($\tau$ $e$)} \\
$\letin{X}{\tau}{\sigma}$ & \code{type $X$ = $\tau$ in $\sigma$} \\
\bottomrule
\end{tabular}
\caption{Concrete syntax of the surface language}
\label{tab:concrete}
\end{table}

Unions have no concrete syntax of their own, and a dictionary with a single key must be written
with at least one guard. The following derived forms, which the parser expands, cover the other
cases. In the expansions, $b$ and $P$ are fresh variables.

\begin{description}
\item[Predicates block.] The block
  \code{predicates $P_1 S_1$ \ldots\ $P_n S_n$ end $\tau$} introduces the predicates
  $P_1, \ldots, P_n$. Each $P_i$ is in scope in the signatures that follow it and in $\tau$:
  \[ \union{(\iabs{P_1}{\pred{S_1}}{\ \cdots\ \union{(\iabs{P_n}{\pred{S_n}}{\tau})}})} \]
\item[Multiple keys.] A dictionary may bind several keys,
  \code{\{ [$a_1$ :\ $c_1$, \ldots, $a_n$ :\ $c_n$] $\overline{\varphi}$ :\ $\tau$ \}}. This nests
  one dictionary for each key. The guards may mention every key, so they belong to the innermost
  dictionary, and the dictionaries around it have no guards:
  \[ \gdict{a_1}{c_1}{}{\ \cdots\ \gdict{a_n}{c_n}{\overline{\varphi}}{\tau}} \]
\item[Unguarded dictionary.] A dictionary written without guards,
  \code{\{ [$a_1$ :\ $c_1$, \ldots, $a_n$ :\ $c_n$] :\ $\tau$ \}}, is guarded by an anonymous
  predicate:
  \[ \union{(\iabs{P}{\pred{(a_1 : c_1, \ldots, a_n : c_n \mid\,)}}{\
       \code{\{ [$a_1$ :\ $c_1$, \ldots, $a_n$ :\ $c_n$] $P(a_1, \ldots, a_n)$ :\ $\tau$ \}}})} \]
\item[String types.] As types, \code{str} stands for $\refine{a}{\str}{\true}$ and \code{"$s$"}
  stands for $\refine{a}{\lit{s}}{\true}$.
\item[Anonymous abstraction.] \code{$q$ => $\tau$} stands for $\iabs{b}{q}{\tau}$.
\item[Repeated arguments and definitions.]
  \code{($\tau$ $r_1$ \ldots\ $r_n$)}, where each $r_i$ is an atom $e$ or \code{!$\sigma$},
  applies $\tau$ to its arguments from left to right. \code{type $X_1$ = $\tau_1$ \ldots\ type $X_n$
  = $\tau_n$ in $\sigma$} nests the definitions, so that each $X_i$ is in scope in the definitions
  that follow it and in $\sigma$.
\end{description}

In a list of dictionary keys, a key written without a sort has sort $\str$.

\FloatBarrier
\section{Surface judgments}
\label{sec:surface-judgments}

The judgments of the surface language are listed below. The rules are syntax-directed, and
subsumption appears only as explicit premises.

\begin{center}
\begin{tabular}{ll}
\toprule
Judgment & Meaning \\
\midrule
$c \leq d$ & string sort $c$ is a subsort of $d$ \\
$\Omega \vdash e : q$ & atom $e$ has sort $q$ \\
$\Omega \proves \varphi$ & $\Omega$ contains a proof of $\varphi$ \\
$\Omega \vdash \varphi \isprop$ & $\varphi$ is a proposition \\
$\Omega \vdash q$ & $q$ is a well-formed sort \\
$\Omega \vdash$ & $\Omega$ is a well-formed sort context \\
$q \leq q'$ & sort $q$ is a subsort of $q'$ \\
$\Omega \vdash \kappa$ & $\kappa$ is a well-formed kind \\
$\kappa \leq \kappa'$ & kind $\kappa$ is a subkind of $\kappa'$ \\
$\discharge{\Omega}{\kappa}{\kappa'}$ & discharging the proofs that $\kappa$ asks for leaves $\kappa'$ \\
$\kinding{\Omega}{\Gamma}{\tau}{\kappa}$ & type $\tau$ has kind $\kappa$ \\
\bottomrule
\end{tabular}
\end{center}

\subsection{Index level}

Figures~\ref{fig:formula-judgments} and~\ref{fig:sort-judgments} give the index-level rules.

\begin{figure}[!htbp]
\judgmentbox{c \leq d}{String subsorting}
\begin{mathpar}
\inferrule*[right=\rulename{C-Str}]{ }{c \leq \str}
\and
\inferrule*[right=\rulename{C-Lit}]{ }{\lit{s} \leq \lit{s}}
\end{mathpar}

\judgmentbox{\Omega \vdash e : q}{Atom sorting}
\begin{mathpar}
\inferrule*[right=\rulename{A-Var}]{(a : q) \in \Omega}{\Omega \vdash a : q}
\and
\inferrule*[right=\rulename{A-Lit}]{ }{\Omega \vdash s : \lit{s}}
\end{mathpar}

\judgmentbox{\Omega \proves \varphi}{Entailment}
\begin{mathpar}
\inferrule*[right=\rulename{E-True}]{ }{\Omega \proves \true}
\and
\inferrule*[right=\rulename{E-Hyp}]{(b : \prf{\alpha}) \in \Omega}{\Omega \proves \alpha}
\end{mathpar}

\judgmentbox{\Omega \vdash \varphi \isprop}{Formula formation}
\begin{mathpar}
\inferrule*[right=\rulename{F-True}]{ }{\Omega \vdash \true \isprop}
\and
\inferrule*[right=\rulename{F-App}]
  {(P : \pred{(a_1 : c_1, \ldots, a_n : c_n \mid \alpha_1, \ldots, \alpha_m)}) \in \Omega \\\\
   \Omega \vdash e_i : d_i \quad d_i \leq c_i \quad (1 \leq i \leq n) \\\\
   \Omega \proves \alpha_j[\overline{e}/\overline{a}] \quad (1 \leq j \leq m)}
  {\Omega \vdash P(e_1, \ldots, e_n) \isprop}
\end{mathpar}

\caption{Atoms and formulas of the surface language}
\label{fig:formula-judgments}
\end{figure}

\begin{figure}[!htbp]
\judgmentbox{\Omega \vdash q}{Sort formation}
\begin{mathpar}
\inferrule*[right=\rulename{S-Str}]{ }{\Omega \vdash c}
\and
\inferrule*[right=\rulename{S-Prf}]{\Omega \vdash \varphi \isprop}{\Omega \vdash \prf{\varphi}}
\and
\inferrule*[right=\rulename{S-Pred}]
  {a_1, \ldots, a_n \text{ distinct} \\\\
   \Omega,\ \overline{a : c},\ b_1 : \prf{\alpha_1},\ \ldots,\ b_{j-1} : \prf{\alpha_{j-1}}
     \vdash \alpha_j \isprop \quad (1 \leq j \leq m)}
  {\Omega \vdash \pred{(a_1 : c_1, \ldots, a_n : c_n \mid \alpha_1, \ldots, \alpha_m)}}
\end{mathpar}

\judgmentbox{\Omega \vdash}{Sort context formation}
\begin{mathpar}
\inferrule*[right=\rulename{Ctx-Empty}]{ }{\emptyctx \vdash}
\and
\inferrule*[right=\rulename{Ctx-Ext}]
  {\Omega \vdash \\ \Omega \vdash q \\ a \notin \mathrm{dom}(\Omega)}
  {\Omega, a : q \vdash}
\end{mathpar}

\judgmentbox{q \leq q'}{Subsorting, which extends string subsorting to all sorts}
\begin{mathpar}
\inferrule*[right=\rulename{Q-Prf}]{ }{\prf{\varphi} \leq \prf{\varphi}}
\and
\inferrule*[right=\rulename{Q-Pred}]
  {c'_i \leq c_i \quad (1 \leq i \leq n) \\
   \{\alpha_1, \ldots, \alpha_m\} = \{\alpha'_1, \ldots, \alpha'_{m'}\}}
  {\pred{(a_1 : c_1, \ldots, a_n : c_n \mid \alpha_1, \ldots, \alpha_m)}
   \leq \pred{(a_1 : c'_1, \ldots, a_n : c'_n \mid \alpha'_1, \ldots, \alpha'_{m'})}}
\end{mathpar}
\caption{Sorts and sort contexts of the surface language}
\label{fig:sort-judgments}
\end{figure}

\paragraph{Atoms.} A variable has the sort that the context assigns to it. A string literal $s$
has the singleton sort $\lit{s}$, which is a subsort of $\str$.

\paragraph{Entailment.} Proofs are irrelevant: all that matters about a proof is that the context
contains one. $\Omega \proves \alpha$ holds when $\Omega$ binds some variable to $\prf{\alpha}$.
Propositions are compared syntactically, so entailment performs no reasoning beyond looking a
proposition up.

\paragraph{Formulas.} Rule \rulename{F-App} is where the two sections of a signature are used. The
string section fixes the number of arguments and their sorts. The proof section is instantiated
with the arguments, and each of the resulting propositions must be entailed by the context. The
arguments must have string sorts, because only string sorts are subsorts of a string sort.

\paragraph{Sorts.} String sorts are always well formed. In rule \rulename{S-Pred}, each proposition
of the proof section is checked under the assumption that the propositions before it hold, so a
proposition may rely on proofs that the same signature asks for earlier. The order of the proof
section matters for this rule only; rule \rulename{Q-Pred} compares proof sections as sets.

\paragraph{Subsorting.} Subsorting on sorts includes the string subsorting rules \rulename{C-Str}
and \rulename{C-Lit}. A predicate may be used where a predicate of a different signature is
expected if it accepts at least the same strings and requires the same proofs. In
\rulename{Q-Pred} the two signatures are first renamed to have the same formals. Proof sorts
are subsorts only of themselves.

\FloatBarrier
\subsection{Kinds}

Figure~\ref{fig:kind-judgments} gives the rules for kinds. A proper type describes a set of
database instances. The kind $\populated$ classifies the proper types whose instances are
guaranteed to hold data; code generation uses it to decide which locations witness that a
predicate holds.

\begin{figure}[!htbp]
\judgmentbox{\Omega \vdash \kappa}{Kind formation}
\begin{mathpar}
\inferrule*[right=\rulename{KF-Proper}]{ }{\Omega \vdash \proper}
\and
\inferrule*[right=\rulename{KF-Populated}]{ }{\Omega \vdash \populated}
\and
\inferrule*[right=\rulename{KF-Arrow}]
  {\Omega \vdash \kappa \\ \Omega \vdash \rho}
  {\Omega \vdash \kappa \to \rho}
\and
\inferrule*[right=\rulename{KF-All}]
  {\Omega \vdash q \\ \Omega, a : q \vdash \kappa}
  {\Omega \vdash \kall{a}{q}{\kappa}}
\end{mathpar}

\judgmentbox{\kappa \leq \kappa'}{Subkinding}
\begin{mathpar}
\inferrule*[right=\rulename{SK-Proper}]{ }{\proper \leq \proper}
\and
\inferrule*[right=\rulename{SK-Populated}]{ }{\populated \leq \populated}
\and
\inferrule*[right=\rulename{SK-Forget}]{ }{\populated \leq \proper}
\and
\inferrule*[right=\rulename{SK-Arrow}]
  {\kappa_2 \leq \kappa_1 \\ \rho_1 \leq \rho_2}
  {\kappa_1 \to \rho_1 \leq \kappa_2 \to \rho_2}
\and
\inferrule*[right=\rulename{SK-All}]
  {q_2 \leq q_1 \\ \kappa_1 \leq \kappa_2}
  {\kall{a}{q_1}{\kappa_1} \leq \kall{a}{q_2}{\kappa_2}}
\end{mathpar}

\judgmentbox{\discharge{\Omega}{\kappa}{\kappa'}}{Proof discharge}
\begin{mathpar}
\inferrule*[right=\rulename{D-Prf}]
  {\Omega \proves \varphi \\ \discharge{\Omega}{\kappa}{\kappa'}}
  {\discharge{\Omega}{\kall{b}{\prf{\varphi}}{\kappa}}{\kappa'}}
\and
\inferrule*[right=\rulename{D-Done}]
  {\kappa \text{ is not of the form } \kall{b}{\prf{\varphi}}{\rho}}
  {\discharge{\Omega}{\kappa}{\kappa}}
\end{mathpar}
\caption{Kind-level judgments of the surface language}
\label{fig:kind-judgments}
\end{figure}

Proof discharge is how proof arguments are supplied without being written.
$\discharge{\Omega}{\kappa}{\kappa'}$ strips the leading proof parameters from $\kappa$, provided
that $\Omega$ entails each of them. If $\kappa$ begins with a proof parameter that $\Omega$ does not
entail, no rule applies. Because proofs are irrelevant, the variable $b$ of a proof parameter does
not occur in the rest of the kind, so \rulename{D-Prf} has nothing to substitute.

\subsection{Kinding}

Figure~\ref{fig:kinding} gives the kinding rules. A \emph{schema} is a type $\tau$ such that
$\kinding{\emptyctx}{\emptyctx}{\tau}{\kappa}$ for some kind $\kappa \leq \proper$.

\begin{figure}[!htbp]
\judgmentbox{\kinding{\Omega}{\Gamma}{\tau}{\kappa}}{Kinding}
\begin{mathpar}
\inferrule*[right=\rulename{K-Var}]
  {(X : \kappa) \in \Gamma}
  {\kinding{\Omega}{\Gamma}{X}{\kappa}}
\and
\inferrule*[right=\rulename{K-Dict}]
  {\Omega,\ a : c,\ b_1 : \prf{\varphi_1},\ \ldots,\ b_{j-1} : \prf{\varphi_{j-1}}
     \vdash \varphi_j \isprop \quad (1 \leq j \leq k) \\\\
   \kinding{\Omega,\ a : c,\ \overline{b : \prf{\varphi}}}{\Gamma}{\tau}{\kappa} \\ \kappa \leq \proper}
  {\kinding{\Omega}{\Gamma}{\gdict{a}{c}{\varphi_1, \ldots, \varphi_k}{\tau}}{\proper}}
\and
\inferrule*[right=\rulename{K-Record}]
  {\kinding{\Omega}{\Gamma}{\tau_i}{\kappa_i} \quad \kappa_i \leq \proper \quad (1 \leq i \leq n)}
  {\kinding{\Omega}{\Gamma}{\{\ell_1 : \tau_1, \ldots, \ell_n : \tau_n\}}
    {\textstyle\bigsqcup_i \kappa_i}}
\and
\inferrule*[right=\rulename{K-Refine}]
  {\Omega, a : c \vdash \varphi \isprop}
  {\kinding{\Omega}{\Gamma}{\refine{a}{c}{\varphi}}{\populated}}
\and
\inferrule*[right=\rulename{K-IndAbs}]
  {\Omega \vdash q \\ \kinding{\Omega, a : q}{\Gamma}{\tau}{\kappa}}
  {\kinding{\Omega}{\Gamma}{\iabs{a}{q}{\tau}}{\kall{a}{q}{\kappa}}}
\and
\inferrule*[right=\rulename{K-TyAbs}]
  {\Omega \vdash \kappa \\ \kinding{\Omega}{\Gamma, X : \kappa}{\tau}{\rho}}
  {\kinding{\Omega}{\Gamma}{\tabs{X}{\kappa}{\tau}}{\kappa \to \rho}}
\and
\inferrule*[right=\rulename{K-Union}]
  {\kinding{\Omega}{\Gamma}{\tau}{\kall{a}{q}{\kappa}} \\ \kappa \leq \proper}
  {\kinding{\Omega}{\Gamma}{\union{\tau}}{\proper}}
\and
\inferrule*[right=\rulename{K-TyApp}]
  {\kinding{\Omega}{\Gamma}{\tau}{\kappa \to \rho} \\
   \kinding{\Omega}{\Gamma}{\sigma}{\kappa'} \\ \kappa' \leq \kappa}
  {\kinding{\Omega}{\Gamma}{\tau\;\sigma}{\rho}}
\and
\inferrule*[right=\rulename{K-IndApp}]
  {\kinding{\Omega}{\Gamma}{\tau}{\kall{a}{q}{\kappa}} \\
   \Omega \vdash e : q' \\ q' \leq q \\
   \discharge{\Omega}{\kappa[e/a]}{\kappa'}}
  {\kinding{\Omega}{\Gamma}{\iapp{\tau}{e}}{\kappa'}}
\and
\inferrule*[right=\rulename{K-Let}]
  {\kinding{\Omega}{\Gamma}{\tau}{\kappa} \\ \kinding{\Omega}{\Gamma, X : \kappa}{\sigma}{\rho}}
  {\kinding{\Omega}{\Gamma}{\letin{X}{\tau}{\sigma}}{\rho}}
\end{mathpar}
In \rulename{K-Record}, $\bigsqcup_i \kappa_i$ is $\populated$ if $\kappa_i = \populated$ for some
$i$, and $\proper$ otherwise.
\caption{Kinding of surface types}
\label{fig:kinding}
\end{figure}

A dictionary has kind $\proper$ even when its values are populated, because it may have no
entries. A record is populated if at least one of its fields is. A refinement is a location that
holds a string, so it is always populated.

Rule \rulename{K-IndApp} applies a type to one atom and then discharges the proof parameters that
follow. For example, suppose $\tau$ has kind
$\kall{a}{\str}{\kall{b}{\prf{P(a)}}{\proper}}$ and $\Omega$ binds $x : \str$. Then $\iapp{\tau}{x}$
has kind $\proper$ if $\Omega \proves P(x)$, and has no kind otherwise. Proof parameters are
discharged only at this point. A type whose kind begins with a proof parameter, such as
$\iabs{b}{\prf{\varphi}}{\tau}$, cannot be applied to anything except a variable $b'$ that $\Omega$
binds to $\prf{\varphi}$. Only a variable can be passed where a predicate or a proof is expected,
because a string literal has a string sort.

Guards put proofs into the context. Rule \rulename{K-Dict} checks each guard under the assumption
that the guards before it hold, and checks the type of the values under the assumption that all of
the guards hold. By the expansion in Section~\ref{sec:concrete}, a predicates block is checked
under \rulename{K-Union} and \rulename{K-IndAbs}, which extend the context with its predicates.

\begin{remark}
The rules do not have $\Omega \vdash$ as a premise. Starting from the empty context, a derivation
extends the context only in \rulename{K-Refine}, which adds a string sort, in \rulename{K-IndAbs}
and \rulename{KF-All}, which check the sort they add, and in \rulename{K-Dict} and
\rulename{S-Pred}, which add string sorts and the propositions they have just checked.
\end{remark}

\FloatBarrier
\section{Core syntax}
\label{sec:core-syntax}

Figure~\ref{fig:core-syntax} gives the abstract syntax of the core language. We use the same
metavariables for the kinds, types, and contexts of both languages; the index and sort levels use
different ones. The core language has no concrete syntax.

\begin{figure}[!htbp]
\[
\begin{array}{llcll}
\text{index} & i, j, k & ::= & a & \text{index variable} \\
  & & \mid & s & \text{string literal} \\
  & & \mid & \true & \text{true proposition} \\
  & & \mid & i\;j & \text{index application} \\[4pt]
\text{sort} & p, r & ::= & \str & \text{string sort} \\
  & & \mid & \lit{s} & \text{literal sort} \\
  & & \mid & \prop & \text{proposition sort} \\
  & & \mid & \prf{i} & \text{proof sort} \\
  & & \mid & (a : p) \to r & \text{function sort} \\[4pt]
\text{kind} & \kappa, \rho & ::= & \proper \;\mid\; \populated \;\mid\; \kappa \to \rho
  \;\mid\; \kall{a}{p}{\kappa} & \\[4pt]
\text{type} & \tau, \sigma & ::= & X & \text{type variable} \\
  & & \mid & \dict{a}{p}{\tau} & \text{dictionary} \\
  & & \mid & \{\ell_1 : \tau_1, \ldots, \ell_n : \tau_n\} & \text{record} \\
  & & \mid & \refine{a}{p}{i} & \text{string refinement} \\
  & & \mid & \iabs{a}{p}{\tau} & \text{index-to-type abstraction} \\
  & & \mid & \tabs{X}{\kappa}{\tau} & \text{type-to-type abstraction} \\
  & & \mid & \union{\tau} & \text{union} \\
  & & \mid & \tau\;\sigma & \text{type application} \\
  & & \mid & \iapp{\tau}{i} & \text{index application} \\
  & & \mid & \letin{X}{\tau}{\sigma} & \text{type definition} \\[4pt]
\text{sort context} & \Omega, \Psi & ::= & \emptyctx \;\mid\; \Omega, a : p & \\[2pt]
\text{kind context} & \Gamma & ::= & \emptyctx \;\mid\; \Gamma, X : \kappa &
\end{array}
\]
\caption{Abstract syntax of the core language}
\label{fig:core-syntax}
\end{figure}

Core indices form a single syntactic category that contains strings, propositions, and functions.
A predicate is an index of function sort, and it is applied to one argument at a time. In the
function sort $(a : p) \to r$ the variable $a$ is bound in $r$; we write $p \to r$ when $a$ does not
occur in $r$. Unlike the surface language, the core language has a sort $\prop$ of propositions,
which is the codomain of predicates, and the sort of a predicate may place its proof arguments
anywhere among its string arguments.

Core types have the same forms as surface types, with the following differences. A core dictionary
has no guards. The key of a dictionary and the variable of a refinement may be annotated with any
sort. The formula of a refinement is an arbitrary index, and a type may be applied to an arbitrary
index.

\FloatBarrier
\section{Translation}
\label{sec:translation}

Figure~\ref{fig:translation} defines the translation $\sem{-}$ from the surface language to the
core language. It is a homomorphism on kinds, and on types other than dictionaries. It does three
things.

\begin{figure}[!htbp]
\textbf{Atoms and formulas}
\begin{align*}
\sem{a} &= a &
\sem{s} &= s &
\sem{\true} &= \true &
\sem{P(e_1, \ldots, e_n)} &= P\;\sem{e_1}\;\cdots\;\sem{e_n}
\end{align*}
\textbf{Sorts}
\begin{align*}
\sem{\str} &= \str &
\sem{\lit{s}} &= \lit{s} &
\sem{\prf{\varphi}} &= \prf{\sem{\varphi}}
\end{align*}
\begin{align*}
\sem{\pred{(a_1 : c_1, \ldots, a_n : c_n \mid \alpha_1, \ldots, \alpha_m)}}
  &= (a_1 : \sem{c_1}) \to \cdots \to (a_n : \sem{c_n}) \to {} \\
  &\phantom{{}={}} \prf{\sem{\alpha_1}} \to \cdots \to \prf{\sem{\alpha_m}} \to \prop
\end{align*}
\textbf{Kinds}
\begin{align*}
\sem{\proper} &= \proper &
\sem{\populated} &= \populated &
\sem{\kappa \to \rho} &= \sem{\kappa} \to \sem{\rho} &
\sem{\kall{a}{q}{\kappa}} &= \kall{a}{\sem{q}}{\sem{\kappa}}
\end{align*}
\textbf{Types}
\begin{align*}
\sem{\gdict{a}{c}{\varphi_1, \ldots, \varphi_k}{\tau}}
  &= \dict{a}{\sem{c}}{\ \union{(\iabs{b_1}{\prf{\sem{\varphi_1}}}{\ \cdots\
       \union{(\iabs{b_k}{\prf{\sem{\varphi_k}}}{\sem{\tau}})}})}}
\end{align*}
\begin{align*}
\sem{X} &= X &
\sem{\{\overline{\ell : \tau}\}} &= \{\overline{\ell : \sem{\tau}}\} \\
\sem{\refine{a}{c}{\varphi}} &= \refine{a}{\sem{c}}{\sem{\varphi}} &
\sem{\iabs{a}{q}{\tau}} &= \iabs{a}{\sem{q}}{\sem{\tau}} \\
\sem{\tabs{X}{\kappa}{\tau}} &= \tabs{X}{\sem{\kappa}}{\sem{\tau}} &
\sem{\union{\tau}} &= \union{\sem{\tau}} \\
\sem{\tau\;\sigma} &= \sem{\tau}\;\sem{\sigma} &
\sem{\iapp{\tau}{e}} &= \iapp{\sem{\tau}}{\sem{e}} \\
\sem{\letin{X}{\tau}{\sigma}} &= \letin{X}{\sem{\tau}}{\sem{\sigma}}
\end{align*}
\textbf{Contexts}
\begin{align*}
\sem{\emptyctx} &= \emptyctx &
\sem{\Omega, a : q} &= \sem{\Omega}, a : \sem{q} &
\sem{\Gamma, X : \kappa} &= \sem{\Gamma}, X : \sem{\kappa}
\end{align*}
In the translation of a dictionary, $b_1, \ldots, b_k$ are fresh variables.
\caption{Translation from the surface language to the core language}
\label{fig:translation}
\end{figure}

\paragraph{Predicate applications are curried.} $P(e_1, \ldots, e_n)$ becomes the left-nested
application $(\cdots(P\;\sem{e_1})\cdots)\;\sem{e_n}$. A nullary application $P()$ becomes the
variable $P$.

\paragraph{Signatures become function sorts.} The string section becomes a sequence of dependent
function sorts, one for each string formal, in order. The proof section becomes a sequence of
function sorts from proofs, in order, placed after all of the string formals. The codomain is
$\prop$. Because the string formals are bound in the proof section of a signature, they are bound
in the translated proof sorts as well. The proof arguments bind no variable.

For example, the signature of $\mathit{CardId}$ from Section~\ref{sec:surface-syntax} translates
as
\begin{align*}
& \sem{\pred{(\mathit{cust} : \str,\ \mathit{card} : \str \mid \mathit{CustId}(\mathit{cust}))}} \\
& \qquad = (\mathit{cust} : \str) \to (\mathit{card} : \str) \to
       \prf{(\mathit{CustId}\;\mathit{cust})} \to \prop
\end{align*}
and the application $\mathit{CardId}(c, x)$ translates to $(\mathit{CardId}\;c)\;x$. The translated
application is not yet a proposition: it is a function that awaits a proof of
$\mathit{CustId}\;c$. The surface rule \rulename{F-App} requires the context to entail
$\mathit{CustId}(c)$, which is what allows the core language to supply that proof.

\paragraph{Guards become unions.} A core dictionary has no guards. Each guard of a surface
dictionary becomes a union over the proofs of the guard, wrapped around the translated type of the
values, in order. Under the abstraction $\iabs{b_j}{\prf{\sem{\varphi_j}}}{\cdots}$ a proof of the
guard is in the context, which is how the core language expresses that the values of the
dictionary may rely on the guard.

\paragraph{Intended property.} The translation is meant to preserve kinding: if
$\kinding{\Omega}{\Gamma}{\tau}{\kappa}$ in the surface language, then
$\kinding{\sem{\Omega}}{\sem{\Gamma}}{\sem{\tau}}{\sem{\kappa}}$ should be derivable in the core
language. This document does not define the core judgments, so the property is not stated
precisely or proven here.

\FloatBarrier
\section{MiniMUMPS}
\label{sec:minimumps}

MiniMUMPS is a small calculus whose expressions read and traverse a database. A schema is compiled
to a MiniMUMPS expression that tests whether a database conforms to it.

Figure~\ref{fig:minimumps} gives the syntax and the reduction relation of MiniMUMPS. We write
$M, N$ for expressions, $V$ for values, and $x$ for variables. Lists are written
$[M_1, \ldots, M_n]$; a list of strings denotes a location in the database.

\begin{figure}[!htbp]
\[
\begin{array}{llcl}
\text{expression} & M, N & ::= & x \;\mid\; s \;\mid\; \true \;\mid\; \false
   \;\mid\; [M_1, \ldots, M_n] \;\mid\; M = N \;\mid\; M \wedge N \\
  & & \mid & \ifte{M}{N_1}{N_2} \;\mid\; \lambda x.\,M \;\mid\; M\;N \\
  & & \mid & \isdef{M} \;\mid\; \readm{M} \;\mid\; \firstten{M}{x}{N} \\[4pt]
\text{value} & V & ::= & s \;\mid\; \true \;\mid\; \false \;\mid\; [V_1, \ldots, V_n]
   \;\mid\; \lambda x.\,M \\[4pt]
\text{evaluation context} & E & ::= & \square \;\mid\; [\overline{V}, E, \overline{M}]
   \;\mid\; E = M \;\mid\; V = E \;\mid\; E \wedge M \\
  & & \mid & \ifte{E}{N_1}{N_2} \;\mid\; E\;M \;\mid\; V\;E \\
  & & \mid & \isdef{E} \;\mid\; \readm{E} \;\mid\; \firstten{E}{x}{N}
\end{array}
\]

\judgmentbox{I \mid M \hookrightarrow M'}{Reduction with read access to $I$}
\begin{mathpar}
\inferrule*[right=\rulename{R-Ctx}]
  {I \mid M \hookrightarrow M'}
  {I \mid E[M] \hookrightarrow E[M']}
\and
\inferrule*[right=\rulename{R-Beta}]{ }{I \mid (\lambda x.\,M)\;V \hookrightarrow M[V/x]}
\and
\inferrule*[right=\rulename{R-IfTrue}]{ }{I \mid \ifte{\true}{N_1}{N_2} \hookrightarrow N_1}
\and
\inferrule*[right=\rulename{R-IfFalse}]{ }{I \mid \ifte{\false}{N_1}{N_2} \hookrightarrow N_2}
\and
\inferrule*[right=\rulename{R-AndTrue}]{ }{I \mid \true \wedge M \hookrightarrow M}
\and
\inferrule*[right=\rulename{R-AndFalse}]{ }{I \mid \false \wedge M \hookrightarrow \false}
\and
\inferrule*[right=\rulename{R-EqTrue}]{ }{I \mid s = s \hookrightarrow \true}
\and
\inferrule*[right=\rulename{R-EqFalse}]{s \neq t}{I \mid s = t \hookrightarrow \false}
\and
\inferrule*[right=\rulename{R-Defined}]
  {I \text{ is active at } [s_1, \ldots, s_n]}
  {I \mid \isdef{[s_1, \ldots, s_n]} \hookrightarrow \true}
\and
\inferrule*[right=\rulename{R-Undefined}]
  {I \text{ is not active at } [s_1, \ldots, s_n]}
  {I \mid \isdef{[s_1, \ldots, s_n]} \hookrightarrow \false}
\and
\inferrule*[right=\rulename{R-Read}]
  {I([s_1, \ldots, s_n]) = s}
  {I \mid \readm{[s_1, \ldots, s_n]} \hookrightarrow s}
\and
\inferrule*[right=\rulename{R-First}]
  {t_1 < \cdots < t_k \text{ are the keys of } I \text{ at } [s_1, \ldots, s_n] \\ m = \min(k, 10)}
  {I \mid \firstten{[s_1, \ldots, s_n]}{x}{M} \hookrightarrow
     M[t_1/x] \wedge \cdots \wedge M[t_m/x] \wedge \true}
\end{mathpar}
\caption{MiniMUMPS}
\label{fig:minimumps}
\end{figure}

\paragraph{Database instances.}
A \emph{database instance} $I$ is a finite partial function from lists of strings to strings. We
write $l \cdot l'$ for the concatenation of the lists $l$ and $l'$. $I$ is \emph{active} at $l$ if
$I(l \cdot l')$ is defined for some list $l'$, which may be empty. The \emph{keys} of $I$ at $l$
are the strings $t$ such that $I$ is active at $l \cdot [t]$. We assume a fixed total order on
strings.

Three forms give access to the database. $\isdef{M}$ tests whether the instance is active at a
location. $\readm{M}$ returns the string stored at a location, and is stuck if there is none.
$\firstten{M}{x}{N}$ evaluates $N$ once for each of the first ten keys at a location, with $x$
bound to the key, and takes the conjunction of the results. Conjunction evaluates its left
operand first and does not evaluate its right operand if the left one is $\false$.

We write $I \mid M \Downarrow V$ when $M$ reduces to the value $V$ in zero or more steps.

\FloatBarrier
\section{Compilation to MiniMUMPS}
\label{sec:compilation}

This section compiles a schema directly to a MiniMUMPS expression, called a \emph{tester}, without
going through the core language. The tester is meant to evaluate to $\true$ on every database
instance that conforms to the schema. It is a testing tool rather than a validator: it can show
that an instance does not conform, but not that it does.

\subsection{Normal types}

Compilation is defined on schemas in normal form. A closed proper type in normal form contains no
type variables, applications, or definitions, and every union in it is taken over an abstraction.
Furthermore, the concrete syntax only produces unions over predicates. The types that we compile
are therefore
\[
\begin{array}{lcl}
\theta & ::= & \gdict{a}{c}{\varphi_1, \ldots, \varphi_k}{\theta}
   \;\mid\; \{\ell_1 : \theta_1, \ldots, \ell_n : \theta_n\} \\
 & \mid & \refine{a}{c}{\varphi}
   \;\mid\; \union{(\iabs{P}{\pred{S}}{\theta})}
\end{array}
\]
The implementation brings a schema into this form by evaluating it before compiling it. Evaluation
is not defined in this document.

\subsection{Locations and decision procedures}

A \emph{location} $\Delta$ is a list of atoms. It describes where in the database the instances of
a type are stored. A location is extended by a dictionary's key variable or by a record's label,
which we treat as a string literal:
\[ \Delta \;::=\; \epsilon \;\mid\; \Delta \cdot e \]
Atoms are MiniMUMPS expressions, so a location $\Delta = e_1 \cdots e_n$ determines the list
expression $[\Delta] = [e_1, \ldots, e_n]$.

We use the following notation for the sections of a sort context:
\begin{align*}
\mathrm{str}(\Omega) &= \{ a \mid (a : c) \in \Omega \} &
\mathrm{prf}(\Omega) &= \{ \alpha \mid (b : \prf{\alpha}) \in \Omega \}
\end{align*}
In the contexts that arise during compilation, $\mathrm{str}(\Omega)$ contains exactly the key
variables of the enclosing dictionaries, all of which occur in the current location.

The compiled test for a refinement $\refine{a}{c}{P(\overline{e})}$ must decide whether
$P(\overline{e})$ holds. The predicate $P$ is abstract, so the test cannot evaluate it. Instead it
relies on the dictionaries that $P$ guards: the keys of a dictionary guarded by $P$ are exactly the
strings that satisfy $P$, provided that every key of the dictionary is physically present.

A \emph{decision procedure} $\dproc{\Psi}{\Delta}$ pairs a sort context with a location. It records
that a guarded dictionary whose values are populated stores its values at $\Delta$, and that
$\Psi$ is the context of those values. The intent is that, for a conforming instance and any choice
of strings for $\mathrm{str}(\Psi)$, the instance is active at $\Delta$ exactly when the
propositions $\mathrm{prf}(\Psi)$ hold.

\begin{figure}[!htbp]
\judgmentbox{\supplies{\Omega}{\Delta}{\theta}{\mathcal{D}}}{$\theta$ supplies the decision procedures $\mathcal{D}$}
\begin{mathpar}
\inferrule*[right=\rulename{U-Refine}]{ }
  {\supplies{\Omega}{\Delta}{\refine{a}{c}{\varphi}}{\emptyset}}
\and
\inferrule*[right=\rulename{U-Record}]
  {\supplies{\Omega}{\Delta \cdot \ell_i}{\theta_i}{\mathcal{D}_i} \quad (1 \leq i \leq n)}
  {\supplies{\Omega}{\Delta}{\{\ell_1 : \theta_1, \ldots, \ell_n : \theta_n\}}
    {\mathcal{D}_1 \cup \cdots \cup \mathcal{D}_n}}
\and
\inferrule*[right=\rulename{U-Pred}]
  {\supplies{\Omega, P : \pred{S}}{\Delta}{\theta}{\mathcal{D}}}
  {\supplies{\Omega}{\Delta}{\union{(\iabs{P}{\pred{S}}{\theta})}}{\mathcal{D}}}
\and
\inferrule*[right=\rulename{U-Dict-Populated}]
  {\Omega' = \Omega,\ a : c,\ \overline{b : \prf{\varphi}} \\ k \geq 1 \\\\
   \supplies{\Omega'}{\Delta \cdot a}{\theta}{\mathcal{D}} \\
   \kinding{\Omega'}{\emptyctx}{\theta}{\populated}}
  {\supplies{\Omega}{\Delta}{\gdict{a}{c}{\varphi_1, \ldots, \varphi_k}{\theta}}
    {\mathcal{D} \cup \{\dproc{\Omega'}{\Delta \cdot a}\}}}
\and
\inferrule*[right=\rulename{U-Dict}]
  {\Omega' = \Omega,\ a : c,\ \overline{b : \prf{\varphi}} \\
   \supplies{\Omega'}{\Delta \cdot a}{\theta}{\mathcal{D}} \\\\
   k = 0 \ \text{ or }\ \kinding{\Omega'}{\emptyctx}{\theta}{\proper}}
  {\supplies{\Omega}{\Delta}{\gdict{a}{c}{\varphi_1, \ldots, \varphi_k}{\theta}}{\mathcal{D}}}
\end{mathpar}
\caption{The decision procedures that a type supplies}
\label{fig:supplies}
\end{figure}

Figure~\ref{fig:supplies} collects the decision procedures of a type. Only a dictionary supplies
one, and only when it has at least one guard and its values have kind $\populated$
(\rulename{U-Dict-Populated}). If the values of a dictionary are not populated, then a key whose
value holds no data leaves no trace in the database, so the dictionary cannot be used to decide
its guards. The context of the decision procedure contains the dictionary's key and all of its
guards.

\subsection{Matching}

A decision procedure is stated in terms of the key variables of the dictionary that supplies it.
To use it at a refinement, those variables must be instantiated with atoms that are available at
the refinement.

Let $\dproc{\Psi}{\Delta'}$ be a decision procedure with $A = \mathrm{prf}(\Psi)$, let $\Omega$ be
the context of a refinement $\refine{a}{c}{\alpha}$, and let $\sigma$ be a substitution of atoms
for the variables $\mathrm{str}(\Psi)$. We say that $\sigma$ \emph{matches} $\dproc{\Psi}{\Delta'}$
to $\alpha$ in $\Omega$, written
$\matches{\Omega}{\sigma}{\dproc{\Psi}{\Delta'}}{\alpha}$,
if all of the following hold.
\begin{description}
\item[M1.] $\alpha \in A\sigma$ and $A\sigma \subseteq \mathrm{prf}(\Omega) \cup \{\alpha\}$.
\item[M2.] Every variable of $\mathrm{str}(\Psi)$ occurs in $A$.
\item[M3.] If $\Psi = \Psi_1, P : \pred{S}, \Psi_2$ and $P$ occurs in $A$, then $\sigma(x) = x$
  for every $x \in \mathrm{str}(\Psi_1)$.
\end{description}

Condition M1 says that the decision procedure decides $\alpha$, and that everything else it
assumes is already known at the refinement. Under M1, if $\alpha$ holds then all of $A\sigma$
holds, so a conforming instance is active at $\Delta'\sigma$.

Condition M2 ensures that the propositions determine every key on the path to the dictionary. A key
that no proposition mentions could be instantiated with anything, and the resulting location would
say nothing about $\alpha$.

Condition M3 concerns predicates that are introduced under a dictionary, which happens when a
predicates block or a dictionary written without guards occurs in the value type of a dictionary. Such a
predicate stands for a different relation under each key of the enclosing dictionaries. The
refinement is in the scope of the same predicate, so by the variable convention its context
$\Omega$ begins with $\Psi_1$ as well, and the keys in $\Psi_1$ must be left as they are. When all
predicates are introduced by a block at the top of the schema, M3 holds trivially.

\subsection{Testers}

\begin{figure}[!htbp]
\judgmentbox{\compiles{\Omega}{\Delta}{\mathcal{D}}{\theta}{M}}{$\theta$ at location $\Delta$ compiles to the tester $M$}
\begin{mathpar}
\inferrule*[right=\rulename{C-Dict}]
  {\compiles{\Omega,\ a : c,\ \overline{b : \prf{\varphi}}}{\Delta \cdot a}{\mathcal{D}}{\theta}{M}}
  {\compiles{\Omega}{\Delta}{\mathcal{D}}{\gdict{a}{c}{\varphi_1, \ldots, \varphi_k}{\theta}}
    {\firstten{[\Delta]}{a}{(\sorttest{a}{c} \wedge M)}}}
\and
\inferrule*[right=\rulename{C-Record}]
  {\compiles{\Omega}{\Delta \cdot \ell_i}{\mathcal{D}}{\theta_i}{M_i} \quad (1 \leq i \leq n)}
  {\compiles{\Omega}{\Delta}{\mathcal{D}}{\{\ell_1 : \theta_1, \ldots, \ell_n : \theta_n\}}
    {M_1 \wedge \cdots \wedge M_n}}
\and
\inferrule*[right=\rulename{C-Refine-True}]{ }
  {\compiles{\Omega}{\Delta}{\mathcal{D}}{\refine{a}{c}{\true}}
    {\isdef{[\Delta]} \wedge \sorttest{\readm{[\Delta]}}{c}}}
\and
\inferrule*[right=\rulename{C-Refine-App}]
  {\dproc{\Psi}{\Delta'} \in \mathcal{D} \\
   \matches{\Omega}{\sigma}{\dproc{\Psi}{\Delta'}}{\alpha}}
  {\compiles{\Omega}{\Delta}{\mathcal{D}}{\refine{a}{c}{\alpha}}
    {\isdef{[\Delta]} \wedge \sorttest{\readm{[\Delta]}}{c} \wedge
     \isdef{[\Delta'\sigma][\readm{[\Delta]}/a]}}}
\and
\inferrule*[right=\rulename{C-Pred}]
  {\compiles{\Omega, P : \pred{S}}{\Delta}{\mathcal{D}}{\theta}{M}}
  {\compiles{\Omega}{\Delta}{\mathcal{D}}{\union{(\iabs{P}{\pred{S}}{\theta})}}{M}}

\end{mathpar}
The sort test $\sorttest{M}{c}$ is defined by $\sorttest{M}{\str} = \true$ and
$\sorttest{M}{\lit{s}} = (M = s)$.
\caption{Compilation of normal types to testers}
\label{fig:compilation}
\end{figure}

Figure~\ref{fig:compilation} gives the compilation rules. The judgment
$\compiles{\Omega}{\Delta}{\mathcal{D}}{\theta}{M}$ compiles the type $\theta$, whose instances are
stored at $\Delta$, to the tester $M$. The set $\mathcal{D}$ contains the decision procedures that
refinements may use. The free variables of $M$ are among $\mathrm{str}(\Omega)$; they are the keys
that the enclosing dictionaries are being traversed with.

A normal schema $\theta$ \emph{compiles to} $M$ if
\[ \supplies{\emptyctx}{\epsilon}{\theta}{\mathcal{D}} \qquad\text{and}\qquad
   \compiles{\emptyctx}{\epsilon}{\mathcal{D}}{\theta}{M} \]
for some $\mathcal{D}$. Compilation thus has two phases: the first collects the decision procedures
of the whole schema, and the second generates the tester, in which a refinement may use a decision
procedure supplied anywhere in the schema.

A dictionary's tester visits keys at the dictionary's location, checks each key against the key
sort, and tests the value stored under it (\rulename{C-Dict}). A record's tester is the conjunction
of its fields' testers (\rulename{C-Record}). A refinement's tester checks that the location is
active and that the string stored there belongs to the refinement's sort. If the refinement has a
predicate application $\alpha$, the tester also consults a decision procedure that matches
$\alpha$: it tests whether the instance is active at the decision procedure's location, after
replacing the refinement variable $a$ with the string read from the refinement's own location
(\rulename{C-Refine-App}). A dictionary's guards and a union's predicate are added to the context and leave no mark on the tester; predicates and
proofs have no representation in the database.

Compilation is a relation rather than a function, because \rulename{C-Refine-App} may choose among
several decision procedures and substitutions. A schema does not compile if one of its refinements
has no matching decision procedure.

\paragraph{Example.} Consider the schema
\begin{center}
\begin{tabular}{l}
\code{predicates CardType(t :\ str) end} \\
\code{\{ cardTypes :\ \{ [t :\ str] CardType(t) :\ "*" \},} \\
\code{\ \ favorite :\ \{ z :\ str | CardType(z) \} \}}
\end{tabular}
\end{center}
The value type \code{"*"} of the dictionary is populated, so the dictionary supplies the decision
procedure
\[ \dproc{\ \mathit{CardType} : \pred{(t : \str \mid\,)},\ t : \str,\ b : \prf{\mathit{CardType}(t)}\ }
         {\ \text{\code{"cardTypes"}} \cdot t\ } \]
At the refinement, the context contains no proofs and the proposition to decide is
$\mathit{CardType}(z)$. The substitution $\sigma = [z/t]$ matches: it sends $\mathit{CardType}(t)$
to $\mathit{CardType}(z)$, and $t$ occurs in that proposition. Omitting the sort tests that are
$\true$, the schema compiles to
\begin{align*}
& \big(\firstten{[\text{\code{"cardTypes"}}]}{t}{} \\
& \qquad (\isdef{[\text{\code{"cardTypes"}}, t]} \wedge
          \readm{[\text{\code{"cardTypes"}}, t]} = \text{\code{"*"}})\big) \\
& \wedge\ \big(\isdef{[\text{\code{"favorite"}}]} \wedge
          \isdef{[\text{\code{"cardTypes"}}, \readm{[\text{\code{"favorite"}}]}]}\big)
\end{align*}
The last conjunct checks that the favorite card type is one of the keys of \code{cardTypes}.

\paragraph{Intended property.} If a database instance $I$ conforms to a schema that compiles to
$M$, then $I \mid M \Downarrow \true$ should hold. Conformance is defined by the denotational
semantics of types, which this document does not give, so the property is not stated precisely or
proven here. The converse is not intended. A tester examines at most ten keys of each dictionary,
and it checks that the strings in refinements are keys of the right dictionaries, but not that the
keys of a dictionary are consistent with the rest of the instance.

\FloatBarrier
\section{Correspondence with the implementation}
\label{sec:implementation}

The syntax, the surface judgments, and the translation to the core language follow the F\#
implementation. The surface syntax is defined in \code{Surface/Syntax.fs}, the core syntax in
\code{Core/Syntax.fs}, and the derived forms in \code{Parser.fsy} and \code{ParseHelpers.fs}.
Table~\ref{tab:implementation} lists the function that decides each judgment.

\subsection{Surface judgments and translation}

The implementation departs from Sections~\ref{sec:surface-syntax}
to~\ref{sec:translation} in the following ways.
\begin{itemize}
\item \emph{Bound variables.} The implementation uses names. Where this document assumes that the
  variables of a context are distinct, the implementation renames a bound variable when it is
  about to add a variable that the context already binds (the \code{enter} functions of
  \code{Surface/Syntax.fs}).
\item \emph{Record labels.} The implementation does not check that the labels of a record are
  distinct.
\item \emph{Context formation.} The implementation never checks $\Omega \vdash$. It relies on the
  observation in the remark at the end of Section~\ref{sec:surface-judgments}.
\item \emph{Schemas.} The implementation does not check that the kind of the top-level type is a
  subkind of $\proper$. Code generation fails on a top-level type that is not proper.
\item \emph{Signatures.} The concrete syntax cannot express a signature with an empty string
  section and a nonempty proof section.
\item \emph{Core kind checker.} The implementation also has a kind checker for the core language
  (\code{Core/KindCheck.fs}), which is not described here. It rejects the translations of some
  schemas that the surface rules accept, because it has no counterpart to \rulename{Q-Prf} or
  \rulename{SK-Populated}.
\end{itemize}

\subsection{Compilation to MiniMUMPS}

The implementation generates testers in \code{CodeGen.fs} and matches decision procedures in
\code{DecisionProcedures.fs}. Section~\ref{sec:compilation} specifies what that code is meant to
compute, but the code differs from it in more ways than the rest of the implementation differs
from the rest of this document.
\begin{itemize}
\item \emph{Target language.} The implementation emits MUMPS source text, not MiniMUMPS. Its
  testers visit every key of a dictionary instead of the first ten, and print a message for each
  discrepancy instead of returning a truth value. It expects the top-level type to be a record,
  because MUMPS stores data in named globals.
\item \emph{One pass.} The implementation generates the tester and collects decision procedures in
  a single traversal. A refinement is left as a hole. The hole is filled when the traversal
  leaves the scope of a predicate that the refinement's proposition or context mentions, using
  only the decision procedures collected so far. As a result, whether a schema compiles can depend
  on the order of the fields of a record.
\item \emph{Populated types.} The implementation does not use the kind $\populated$ in
  \rulename{U-Dict-Populated}. It recomputes whether a type is populated during code generation,
  and treats every record as populated.
\item \emph{Matching.} The implementation has no counterpart to condition M3. In place of
  $\alpha \in A\sigma$ in condition M1, it only requires some proposition of $A$ to apply the same
  predicate as $\alpha$, so it can choose a decision procedure whose location does not depend on
  the string being tested.
\item \emph{Refinements.} The implementation does not emit the tests $\isdef{[\Delta]}$ and
  $\sorttest{\readm{[\Delta]}}{c}$ for a refinement, or the sort test for a dictionary key.
\end{itemize}

\begin{table}[!htbp]
\centering
\begin{tabular}{lll}
\toprule
Judgment & Function & Module \\
\midrule
$c \leq d$ & \code{subStrSort} & \code{Surface.SortCheck} \\
$\Omega \vdash e : q$ & \code{sortCheckAtom} & \code{Surface.SortCheck} \\
$\Omega \proves \varphi$ & \code{proves} & \code{Surface.SortCheck} \\
$\Omega \vdash \varphi \isprop$ & \code{checkFormula}, \code{checkPredApp} & \code{Surface.SortCheck} \\
$\Omega \vdash q$ & \code{wfSort} & \code{Surface.SortCheck} \\
$\Omega \vdash$ & \code{wfSortContext} & \code{Surface.SortCheck} \\
$q \leq q'$ & \code{subSort}, \code{subPredSig} & \code{Surface.KindCheck} \\
$\Omega \vdash \kappa$ & \code{wfKind} & \code{Surface.KindCheck} \\
$\kappa \leq \kappa'$ & \code{subKind} & \code{Surface.KindCheck} \\
$\discharge{\Omega}{\kappa}{\kappa'}$ & \code{applyProofs} & \code{Surface.KindCheck} \\
$\kinding{\Omega}{\Gamma}{\tau}{\kappa}$ & \code{kindSynth} & \code{Surface.KindCheck} \\
$\sem{-}$ & \code{compileTy} and its helpers & \code{Surface.Compile} \\
$\supplies{\Omega}{\Delta}{\theta}{\mathcal{D}}$ & \code{genCode}, \code{canonicalize} & \code{CodeGen}, \code{DecisionProcedures} \\
$\matches{\Omega}{\sigma}{\dproc{\Psi}{\Delta'}}{\alpha}$ & \code{instantiate2} & \code{DecisionProcedures} \\
$\compiles{\Omega}{\Delta}{\mathcal{D}}{\theta}{M}$ & \code{genCode} & \code{CodeGen} \\
\bottomrule
\end{tabular}
\caption{Where each judgment is implemented}
\label{tab:implementation}
\end{table}

\end{document}
